\section{Principalization and lifted additive actions}
\label{sec:bundle}

The ideal defining $G$ becomes principal after pullback to a line
bundle with its zero section removed. We construct that bundle and then
show that every homogeneous additive action on $G$ lifts to the
pullback. This replaces the two generators of the ideal by one element
while retaining both degree and fiber weight.

We retain the coefficient ring and ideal from \Cref{prop:graded}:
\[
 R=\kk[x,y,z,u]/\bigl(xy-z(z+1)\bigr),\qquad
 s_R=y-x(x-u^3),\qquad
 I=(f,g)=(x-s_R^2-u^3,\ z-s_Rx).
\]
The inclusion $R\subset G=R[\tau,I/\tau^2]$ is the actual coefficient
inclusion in $R[\tau,\tau^{-1}]$. In terms of the initial generators,
it is given by
\begin{equation}
\label{bun:coefficient-inclusion}
 x=s_R^2+u^3+\tau^2(f/\tau^2),\qquad
 y=s_R+x(x-u^3),\qquad
 z=s_Rx+\tau^2(g/\tau^2).
\end{equation}
This is the inclusion used in the base change below. Its elements have
valuation degree zero.

\subsection{A line bundle that principalizes the center}

The bundle used here is the standard quotient of $\mathrm{SL}_2$ by
its diagonal torus, with the coordinate $u$ adjoined; see
Blanc--van Santen~\cite[p.~8431]{BlancVanSanten} for this
description of the affine quadric. We give its charts explicitly
because the transition functions determine the fiber weights.

Set
\begin{equation}
\label{bun:determinant-ring}
 \Rt=\kk[a,d,b,c,u]/(ac-bd-1),\qquad
 x=ab,\quad y=dc,\quad z=db,
\end{equation}
and let
\begin{equation}
\label{bun:v}
 v=a^2(x-u^3)-d^2=a^3b-a^2u^3-d^2.
\end{equation}
The assignments in \Cref{bun:determinant-ring} define a homomorphism
$R\to\Rt$, since $xy=(ac)(bd)=z(z+1)$.

\begin{proposition}
\label{prop:principalization}
The morphism $\Spec\Rt\to\Spec R$ is the complement of the zero
section in a line bundle, with fiber scaling specified by the weights
\[
 \wt(a)=\wt(d)=1,\qquad
 \wt(b)=\wt(c)=-1,\qquad \wt(u)=0.
\]
The homomorphism $R\to\Rt$ is faithfully flat, and
$I\Rt=(v)$. Moreover, base change gives a faithfully flat inclusion
$G\hookrightarrow\Gt$, where
\begin{equation}
\label{bun:pullback-algebra}
 \Gt:=G\otimes_R\Rt
   =\Rt[\tau,v/\tau^2]\subset\Rt[\tau,\tau^{-1}],\qquad
 \Gt\cong\Rt[\tau,V]/(\tau^2V-v).
\end{equation}
Thus $V$ denotes the element $v/\tau^2$ in the displayed subalgebra.
The ring $\Gt$ is a domain and has the following commuting gradings:
\begin{center}
\begin{tabular}{c@{\qquad}rrrrr}
\toprule
 & $a,d$ & $b,c$ & $u$ & $\tau$ & $V$\\
\midrule
Valuation degree & $0$ & $0$ & $0$ & $-1$ & $2$\\
Fiber weight & $1$ & $-1$ & $0$ & $0$ & $2$\\
\bottomrule
\end{tabular}
\end{center}
\end{proposition}

\begin{proof}
First, $\Rt$ is a domain: the polynomial $ac-(bd+1)$ is primitive
as a polynomial in $c$ over the unique factorization domain
$\kk[a,d,b,u]$, and is irreducible over its fraction field because
it is linear. The element $v$ is nonzero, since $u$ is an
indeterminate over the determinant ring and its coefficient in degree
three in $v$ is $-a^2\ne0$.

The principal opens $D(z+1)$ and $D(z)$ cover $\Spec R$.
Over the first open, $ac=z+1$ is invertible, and the formulas
\begin{equation}
\label{bun:first-chart}
 c=\frac{z+1}{a},\qquad b=\frac{x}{a},\qquad
 d=\frac{ay}{z+1}
\end{equation}
give mutually inverse identifications
\[
 \Rt[(z+1)^{-1}]
   \cong R[(z+1)^{-1}][a,a^{-1}].
\]
Indeed, they give $ac=z+1$, $ab=x$, $dc=y$, and
$db=xy/(z+1)=z$. Over the second open, the analogous formulas are
\begin{equation}
\label{bun:second-chart}
 b=\frac{z}{d},\qquad c=\frac{y}{d},\qquad a=\frac{dx}{z},
 \qquad
 \Rt[z^{-1}]\cong R[z^{-1}][d,d^{-1}].
\end{equation}
Here $db=z$, $dc=y$, $ab=x$, and $ac=xy/z=z+1$.
On the overlap the fiber parameters satisfy
\[
 d=\frac{y}{z+1}a,\qquad
 \frac{y}{z+1}\cdot\frac{x}{z}=1.
\]
Thus their transition function is a base unit. Gluing affine lines
by this same transition function gives a line bundle whose complement
of the zero section is $\Spec\Rt$. The indicated weights are exactly
scalar multiplication on the two fiber parameters. Each Laurent
polynomial chart is faithfully flat over its base chart: as a module
it is free with basis the integral powers of its fiber parameter.
Since the base charts cover, $R\to\Rt$ is faithfully flat.

We prove principalization by global identities. In $\Rt$ put
\[
 h=x-u^3,\qquad C_1=c^2-b^2h,\qquad
 \alpha=a^2(1-2bd),\qquad
 \beta=d^2(3+2bd)+\alpha h.
\]
Since $s_R=dc-abh$ and $ac-bd=1$, we have
\[
 as_R-d=-bv,\qquad
 g=b(d-as_R)=b^2v.
\]
Expansion also gives
\[
 f-C_1v
  =h-(dc-abh)^2-(c^2-b^2h)(a^2h-d^2)
  =h\bigl(1-(ac-bd)^2\bigr)=0.
\]
Finally, with $q=bd$ the coefficient identity is
\begin{align*}
 \alpha C_1+\beta b^2
 &= (ac)^2(1-2bd)+3b^2d^2+2b^3d^3\\
 &= (1+q)^2(1-2q)+3q^2+2q^3=1.
\end{align*}
Consequently
\begin{equation}
\label{bun:principalization-certificate}
 f=C_1v,\qquad g=b^2v,\qquad v=\alpha f+\beta g,
\end{equation}
which proves $I\Rt=(v)$.

Flatness permits us to tensor the injection
$G\hookrightarrow R[\tau,\tau^{-1}]$ with $\Rt$.
The resulting injection has image generated by
$\Rt$, $\tau$, $f/\tau^2$, and $g/\tau^2$.
By \Cref{bun:principalization-certificate}, this image is precisely
$\Rt[\tau,v/\tau^2]$. In particular it is a domain. The map
$G\to G\otimes_R\Rt$ is faithfully flat by base change, and hence
injective.

For completeness, the single displayed relation is the entire
presentation. Let $W$ be an indeterminate and set
$Q=\Rt[\tau,W]/(\tau^2W-v)$.
If $\tau F=(\tau^2W-v)B$ in $\Rt[\tau,W]$, reduction modulo
$\tau$ gives $vB(0,W)=0$. Since $\Rt[W]$ is a domain and $v\ne0$,
we have $B=\tau B_1$; cancellation then gives
$F=(\tau^2W-v)B_1$. Thus $\tau$ is a nonzerodivisor in $Q$.
After inverting $\tau$, the relation eliminates $W$ and identifies
$Q[\tau^{-1}]$ with $\Rt[\tau,\tau^{-1}]$.
It follows that $Q$ injects into that Laurent ring with image
$\Rt[\tau,v/\tau^2]$, proving \Cref{bun:pullback-algebra}.
The two gradings follow from $\deg(v)=0$ and $\wt(v)=2$;
the relation $\tau^2V=v$ is homogeneous for each.
\end{proof}

Write $Y=\Spec R$ and $Z=\Spec G$. The geometric content of
\Cref{bun:pullback-algebra} is the Cartesian square
\[
\begin{tikzcd}[column sep=large,row sep=large]
 \Spec\Gt \arrow[r] \arrow[d] & \Spec\Rt \arrow[d] \\
 Z \arrow[r] & Y.
\end{tikzcd}
\]
If $L_Y$ is the line bundle just constructed on $Y$, then
$\Spec\Gt=L^\times$, where $L$ is its pullback to $Z$ and
$L^\times$ denotes the complement of its zero section.

\subsection{A normalized lift on a line bundle}

The pullback is a line bundle over $Z=\Spec G$. An additive action on
$Z$ need not descend to $Y=\Spec R$, so a lift obtained by first acting
on $Y$ would not suffice. Instead, we prove a lifting statement for any
line bundle over a smooth integral affine base. The two ingredients
are invariance of the Picard group under adjoining a variable and the
absence of nonconstant polynomial units over a domain.

For a variety $X$, $\Pic(X)$ is the group of isomorphism classes of line
bundles on $X$, with tensor product as its operation. For a line bundle
$L$ on $X$, normalization means that the restriction over
$\{0\}\times X$ is the identity on $L$.
The Picard-group assertion is the regular-ring case of
Traverso's theorem~\cite[Theorem~3.6]{Traverso}. Additive
linearization also follows from the equivariant Picard-group results
of Brion~\cite[Proposition~2.10 and Theorem~2.14]{Brion}. For the normalized
isomorphism and cocycle verification we follow
\cite[Lemma~2.9]{Brion}; we include the argument to make the subsequent
compatibility with valuation degree explicit.

\begin{proposition}
\label{prop:linearization}
Let $X$ be a smooth integral affine variety over $\kk$, and let
$\pi\colon\Aff^1\times X\to X$ be the projection. Then
\[
 \pi^*\colon\Pic(X)\xrightarrow{\ \sim\ }
                  \Pic(\Aff^1\times X).
\]
For every algebraic additive action
$\alpha\colon\Ga\times X\to X$ and every line bundle $L$ on $X$,
there is a unique isomorphism of line bundles
\begin{equation}
\label{bun:linearization-map}
 \phi\colon\pi^*L\xrightarrow{\ \sim\ }\alpha^*L
\end{equation}
whose restriction to $\{0\}\times X$ is $\mathrm{id}_L$.
This isomorphism obeys the action cocycle. It therefore defines a
unique action on $L$ that lifts $\alpha$, is linear on its fibers,
and acts as the identity at the group identity. The action preserves
$L^\times$ and commutes with fiber scaling. If $L^\times$ is affine,
the induced derivation of its coordinate ring is locally nilpotent
and preserves the fiber weights.
\end{proposition}

\begin{proof}
Put $C=\Gamma(X,\mathcal O_X)$ and $K=\Frac C$.
Both $X$ and $\Aff^1\times X$ are regular. Their local rings are
factorial, so every Weil divisor is Cartier and their Picard groups
identify with their divisor class groups~\cite[Tags~056S, 0AG0, and~0BE9]{Stacks}.
Let $D$ be a Weil divisor on
$\Spec C[t]$. Its restriction to $\Spec K[t]$ is principal, since
$K[t]$ is a principal ideal domain. Choose $r\in K(t)^*$ giving
that divisor. Subtracting $\operatorname{div}(r)$ from $D$ removes
every prime component whose defining prime ideal has zero
contraction to $C$.

Each remaining component has a height-one prime ideal
$\mathfrak P\subset C[t]$ with nonzero contraction. Choose
$0\ne c\in\mathfrak P\cap C$, and choose a prime
$\mathfrak q$ minimal over $(c)$ and contained in
$\mathfrak P\cap C$. The principal ideal theorem~\cite[Tag~00KV]{Stacks}
and the fact that
$C$ is a domain give $\operatorname{ht}\mathfrak q=1$.
Now $\mathfrak qC[t]$ is a nonzero prime contained in the
height-one prime $\mathfrak P$, so
$\mathfrak P=\mathfrak qC[t]$.
The corresponding divisor is the pullback of the prime divisor
defined by $\mathfrak q$, with multiplicity one: a uniformizer
of the discrete valuation ring $C_{\mathfrak q}$ is also a
uniformizer at the generic point of that pulled-back divisor.
Thus every divisor class on $\Aff^1\times X$ is pulled back
from $X$, proving surjectivity of $\pi^*$. Restriction to
$i_0\colon X\to\Aff^1\times X$, $x\mapsto(0,x)$, is a left
inverse to $\pi^*$ on Picard groups, proving injectivity as well.

Since $\alpha\circ i_0=\pi\circ i_0=\mathrm{id}_X$, the line
bundles $\alpha^*L$ and $\pi^*L$ have the same restriction to
the identity section. The Picard isomorphism just proved implies
that they are isomorphic. Take an isomorphism $\phi_1$ between
them, with the orientation in \Cref{bun:linearization-map}.
Its restriction $i_0^*\phi_1$ is an automorphism of $L$.
Composing on the source with
$\pi^*((i_0^*\phi_1)^{-1})$ produces an isomorphism $\phi$
whose restriction is the identity.

An automorphism of any line bundle on a scheme is multiplication
by a global unit: its endomorphism sheaf is canonically the
structure sheaf. Because $C$ is a domain,
\begin{equation}
\label{bun:parameter-units}
 C[t]^*=C^*,\qquad C[s,t]^*=C^*.
\end{equation}
Indeed, total degree of a product of nonzero polynomials over a
domain is the sum of their total degrees, so a polynomial unit
has degree zero. The ratio of two normalized isomorphisms in
\Cref{bun:linearization-map} is a unit of $C[t]$ whose value at
$t=0$ is one. By \Cref{bun:parameter-units}, that ratio is one.
This proves uniqueness, without a restriction on the units of $C$.

To check the cocycle, denote the fiber map by
$\phi_{t,x}\colon L_x\to L_{\alpha(t,x)}$.
Over $\Aff^2\times X$, the two compositions to compare are
\[
 \phi_{t,\alpha(s,x)}\circ\phi_{s,x},\quad
 \phi_{s+t,x}
 \colon L_x\longrightarrow L_{\alpha(s+t,x)}.
\]
Their sources and targets agree by the action law for $\alpha$.
Thus
\[
 \phi_{s+t,x}^{-1}\circ
 \phi_{t,\alpha(s,x)}\circ\phi_{s,x}
\]
is an automorphism of the pullback of $L$ from $X$, hence a unit
of $C[s,t]$. At $s=0$ it is the identity. By
\Cref{bun:parameter-units}, it is the identity everywhere.
This proves the cocycle. The family of fiber maps consequently
defines an algebraic additive action on the total space of $L$.
Linearity preserves zero and nonzero vectors and commutes with
scalar multiplication, giving the asserted action on $L^\times$.

Finally, suppose $L^\times=\Spec T$ is affine. Write its action
on functions as
\[
 \rho(r)=\sum_{n\ge0}\delta_n(r)t^n\in T[t],\qquad r\in T,
\]
where the sum is finite for each $r$. The identity and action laws
give $\delta_0=\mathrm{id}$ and
\[
 \delta_i\delta_j=\binom{i+j}{i}\delta_{i+j}.
\]
The coefficient of $t$ in multiplicativity of $\rho$ says that
$\delta_1$ is a derivation. Since the characteristic is zero, the
displayed composition identity gives
$\delta_n=\delta_1^n/n!$. Finiteness of $\rho(r)$ now says that
$\delta_1$ is locally nilpotent. Commutation with fiber scaling
says precisely that this derivation preserves the fiber weights.
\end{proof}

\subsection{Compatibility with valuation degree}

The preceding proposition gives a lift preserving fiber weights. We
still need its valuation degree to agree with that of the original
derivation. The degree action has its own explicit bundle lift;
uniqueness of the normalized additive lift makes the two compatible.

\begin{proposition}
\label{prop:homogeneous-lift}
Let $D_0$ be a homogeneous locally nilpotent derivation of $G$
of valuation degree shift $j$. There is a locally nilpotent
derivation $\Dt$ of $\Gt$ restricting to $D_0$ on $G$, with
fiber weight zero and valuation degree shift $j$. If $D_0\ne0$,
then $\Dt\ne0$.
\end{proposition}

\begin{proof}
By \Cref{prop:graded}, $Z=\Spec G$ is smooth and integral.
Exponentiating $D_0$ gives the algebraic action on functions
\[
 \exp(tD_0)(r)=\sum_{n\ge0}\frac{D_0^n(r)}{n!}t^n.
\]
Apply \Cref{prop:linearization} to this action and the line bundle
$L$ on $Z$ whose complement is $\Spec\Gt$.
The resulting locally nilpotent derivation $\Dt$ restricts to
$D_0$ because the action covers the given action on $Z$, and
has fiber weight zero because the lift commutes with fiber
scaling. Since $G\hookrightarrow\Gt$, a nonzero $D_0$ gives a
nonzero $\Dt$.

It remains to establish its valuation degree. For $\lambda\in\kk^*$,
let $\sigma_\lambda$ be the automorphism of $G$ multiplying a
homogeneous element of degree $h$ by $\lambda^h$.
It fixes the coefficient subring $R$ pointwise. Therefore it has
the pullback lift
$\widetilde\sigma_\lambda=\sigma_\lambda\otimes\mathrm{id}_{\Rt}$
on $\Gt$. Explicitly,
\begin{equation}
\label{bun:degree-scaling}
 \widetilde\sigma_\lambda(r)=r\quad(r\in\Rt),\qquad
 \widetilde\sigma_\lambda(\tau)=\lambda^{-1}\tau,\qquad
 \widetilde\sigma_\lambda(V)=\lambda^2V.
\end{equation}
Geometrically this lift acts on $L=Z\times_Y L_Y$ by the degree
action on $Z$ and the identity on $L_Y$. In particular it extends
to a linear bundle automorphism, and its action on the complement
agrees with \Cref{bun:degree-scaling}.

On the coordinate ring of the base, homogeneity of $D_0$ gives
\begin{equation}
\label{bun:base-conjugacy}
 \sigma_\lambda D_0\sigma_\lambda^{-1}=\lambda^jD_0.
\end{equation}
Conjugate the lifted additive action by
$\widetilde\sigma_\lambda$, using the expression on coordinate rings
\[
 \widetilde\sigma_\lambda\,\exp(t\Dt)\,
 \widetilde\sigma_\lambda^{-1}.
\]
This is a normalized linear lift of the base action
$\exp(\lambda^jtD_0)$, by \Cref{bun:base-conjugacy}.
The action $\exp(\lambda^jt\Dt)$ is another normalized linear
lift of that same base action. Uniqueness in
\Cref{prop:linearization} therefore gives equality of these
two families. Taking their coefficients of $t$ yields
\[
 \widetilde\sigma_\lambda\Dt
 \widetilde\sigma_\lambda^{-1}=\lambda^j\Dt
 \qquad(\lambda\in\kk^*).
\]
For a homogeneous $r$ of degree $h$, this says
$\widetilde\sigma_\lambda(\Dt r)=\lambda^{h+j}\Dt r$
for every $\lambda$. The distinct integral characters of
$\Gm$ are linearly independent, so $\Dt r$ is homogeneous of
degree $h+j$ whenever it is nonzero. This proves the required
degree shift.
\end{proof}

The additive lift in \Cref{prop:homogeneous-lift} is constructed
on the line bundle over $Z=\Spec G$. It does not require
$D_0(R)\subset R$, and it does not assert $\Dt(\Rt)\subset\Rt$.
The invariance of coefficient rings used in the proof belongs to
the separate degree action: $\sigma_\lambda$ fixes $R$, and
$\widetilde\sigma_\lambda$ fixes $\Rt$.
